\section{Conclusion}
\label{sec:conclusion}

We studied which notion of flatness governs continual generalization in LoRA-based PECL.
Through a sequential hierarchical PAC-Bayesian analysis, we showed that under a frozen backbone the KL complexity and the sharpness-dependent empirical term in the bound are both supported on the task-admissible subspace $\mathcal{W}_{\Delta,t}$---not on the ambient full-parameter space.
This identifies \emph{admissible-update flatness} as the bound-relevant notion in this regime, while ambient full-space flatness (including the frozen backbone directions) is not what the bound controls.
The analysis yields four testable predictions about perturbation scope and type, all of which are confirmed on vision and language benchmarks across three LoRA organizations and multiple model scales.

In practice, the results support three concrete design principles: restrict perturbations to the adapter scope; prefer gradient-aligned sharpness types (AS$^{(1),\Delta}$); and treat LoRA organization as a structural prior that determines the geometry of $\mathcal{W}_{\Delta,t}$ and hence the effective target for flatness control.
These principles provide a theory-grounded foundation for flatness-aware design in parameter-efficient continual learning.

\section*{Impact Statement}
This paper advances the theoretical and empirical understanding of flatness-aware optimization for parameter-efficient continual learning.
Potential downstream impacts depend on how continual learning systems are deployed; as with other modeling advances, appropriate validation and governance are needed before safety-sensitive applications.
